% Appendix for the Reproducible Science paper: retractions, errors and misconduct in published results.
% Draft v1r. Every count is in ../failure_cases/results/case_counts.json or
% ../failure_cases/results/rw_tag_counts.json; the case table is generated from
% ../failure_cases/data/cases_draft.json, whose quotations are pinned in ../failure_cases/claims/.

\section{Retractions, errors and misconduct in published results}
\label{sec:failure-cases}

Reproducible Science checks that a manuscript agrees with its artifacts, which covers some of the
reasons published results are retracted and not others. To determine which reasons the checks
cover, we examine 34 cases of retraction, correction, or misconduct finding, and classify each
by the cause stated in its primary record.

\paragraph{Detected failure modes.} Table~\ref{tab:fc-modes} gives the coverage of
each failure mode, counting each case once under its principal cause. The checks
cover a stated problem in 14 of the 34 cases: the principal cause in 8 (6 detected and 2
reported as unchecked), and missing data reported
alongside the principal cause in 6 others. The checks detect three failure modes directly: (1)~a printed value that
disagrees with the authors' own records fails the check of that number against its results
file, (2)~a report that departs from its registered plan differs from the frozen and
timestamped plan, and (3)~a reference that does not exist cannot be pinned to a source.

\vspace{0.5em}
\begin{table}[h]
\centering
\small
\caption{Failure modes among the 34 cases and the coverage of each by Reproducible Science. Each case is
counted once, under its principal cause. $^{\dagger}$Added after the records were read.}
\label{tab:fc-modes}
\begin{tabular}{@{}>{\raggedright\arraybackslash}p{0.33\textwidth}r>{\raggedright\arraybackslash}p{0.52\textwidth}@{}}
\toprule
\textbf{Failure mode} & \textbf{Cases} & \textbf{Coverage} \\
\midrule
printed value disagrees with the authors' own records & 3 & \textbf{detected}: the printed number fails the check against its results file \\
report departs from the registered plan & 1 & \textbf{detected}: the report differs from the frozen, timestamped plan \\
references that do not exist & 2 & \textbf{detected}: a reference that does not exist cannot be pinned to a source \\
\midrule
data could not be produced for audit & 2 & \textbf{reported as unchecked}: the sealed digest records that the file existed \\
error inside the analysis & 9 & \textbf{not detected}: the sealed script and inputs allow the analysis to be rerun \\
data fabricated or falsified & 8 & \textbf{not detected} if fabricated before sealing; a later change to a sealed file is detected \\
\midrule
image manipulation or duplication & 4 & \textbf{not addressed} \\
study conducted differently from its description$^{\dagger}$ & 4 & \textbf{not addressed} \\
none of the above & 1 & \textbf{not addressed} \\
\midrule
Total & 34 & \\
\bottomrule
\end{tabular}
\end{table}

\paragraph{Undetected failure modes.} Data that cannot be produced for an audit are reported as
unchecked rather than verified. An error inside the analysis passes every check, because the
manuscript agrees with the output of the faulty analysis, but the sealed script and inputs
allow a reader to rerun the analysis and locate the error. Reproducible Science does not address data
fabricated before sealing, manipulated images, or a study conducted differently from its
published description.

\paragraph{Case selection.} The primary record of a case is its retraction notice, a finding of
the United States Office of Research Integrity, or a published reanalysis. The cases comprise
nine of the ten most-cited retracted papers listed by Retraction Watch
\citep{retractionwatch2026leaderboard}, publications named in Office of Research
Integrity findings, retractions drawn from the Retraction Watch dataset
\citep{crossref2026retractionwatchdata} by reason tag, and well-known examples such as
the Duke chemotherapy predictors and the Surgisphere studies. We select the cases as examples
of each failure mode. We define ten failure modes before reading any record and add an
eleventh for four records that fit none of the ten.

\paragraph{Retraction Watch dataset.} To place the cases in context, we analyze the reasons
recorded in the Retraction Watch dataset, which covers 67,530 retractions. We exclude the
24,166 retractions attributed to paper mills or compromised peer review, because these concern
the publication process, and group the reason tags of the remaining 43,364
(Table~\ref{tab:fc-dataset}). Among these, 4,273 retractions (9.8\%) cite an error in the
analysis, data, methods, results or text, and 1,909 (4.4\%) cite original data or images that
were not provided. Retractions that cite unavailable data or images follow publication by a
median of 4.91 years, compared with 1.31 years for retractions that cite an error.

\vspace{0.5em}
\begin{table}[h]
\centering
\small
\caption{Retractions in the Retraction Watch dataset by group of reason tags, among the 43,364
not tagged for paper mills or compromised peer review. A retraction can fall in several groups.
Years run from the paper's date to the retraction's date.}
\label{tab:fc-dataset}
\begin{tabular}{@{}>{\raggedright\arraybackslash}p{0.44\textwidth}rrr@{}}
\toprule
\textbf{Reason for retraction} & \textbf{Retractions} & \textbf{Share} & \textbf{Median years to retraction} \\
\midrule
image duplication, manipulation, fabrication or concerns & 7,171 & 16.5\% & 4.39 \\
error stated in analysis, data, methods, results or text & 4,273 & 9.8\% & 1.31 \\
original data or images not provided & 1,909 & 4.4\% & 4.91 \\
falsification or fabrication of data & 1,659 & 3.8\% & 4.27 \\
results not reproducible & 861 & 2.0\% & 2.01 \\
limited or no information in the notice & 11,957 & 27.6\% & 0.11 \\
\bottomrule
\end{tabular}
\end{table}

\paragraph{Cases.} Table~\ref{tab:fc-cases} gives the cause quoted from the primary record of
each case.

{\small
\begin{longtable}{@{}>{\raggedright\arraybackslash}p{0.27\textwidth}>{\raggedright\arraybackslash}p{0.68\textwidth}@{}}
\caption{Thirty-four cases, grouped by failure mode, with the cause in the words of the primary record: a retraction notice, a finding of the Office of Research Integrity, or a published reanalysis.}
\label{tab:fc-cases} \\
\toprule
\textbf{Case} & \textbf{What its primary record says} \\
\midrule
\endfirsthead
\toprule
\textbf{Case} & \textbf{What its primary record says} \\
\midrule
\endhead
\bottomrule
\endfoot
\multicolumn{2}{@{}l}{\textbf{Printed value disagrees with the authors' own records}} \\
Abegaz et al., PLOS ONE, 2014 & ``the estimated number of deaths and injuries reported cannot be reproduced when using the unmatched and matched police and hospital data'' \\
Sun et al., PLOS ONE, 2024 & ``the p50 median value is lower than the reported minimum value'' Also states that data, images or records could not be produced. \\
Wansink et al., Archives of Pediatrics and Adolescent Medicine, 2012 & ``We had erroneously reported the age group as children ranging from 8 to 11 years old; however, the children were 3 to 5 years old.'' \\
\midrule
\multicolumn{2}{@{}l}{\textbf{Report departs from the registered plan}} \\
67 trials in five medical journals, 2015 to 2016 (a cohort, not one publication) & ``365 novel outcomes were reported without declaration in total, a mean of 5.4 per trial'' \\
\midrule
\multicolumn{2}{@{}l}{\textbf{References that do not exist}} \\
Khalilzadeh et al., PLOS ONE, 2026 & ``Citation of unretrievable references 1, 2, 4, 5, 8, 10, 12--18, 20--22, 33, 36, 51, 61, 75, and 79.'' \\
Mohammad et al., PLOS ONE, 2026 & ``Citation of unretrievable References 2, 3, 6, 7, 8, 18, 20, and 26.'' \\
\midrule
\multicolumn{2}{@{}l}{\textbf{Data could not be produced for audit}} \\
Mehra et al., Lancet, 2020 & ``Surgisphere would not transfer the full dataset, client contracts, and the full ISO audit report to their servers for analysis'' \\
Mehra et al., New England Journal of Medicine, 2020 & ``the raw data could not be made available to a third-party auditor, we are unable to validate the primary data sources'' \\
\midrule
\multicolumn{2}{@{}l}{\textbf{Error inside the analysis}} \\
Morisky et al., Journal of Clinical Hypertension, 2008 & an independent statistical review ``concluded that the results were misleading due to issues regarding the sensitivity and specificity of the medical adherence scale used'' \\
Gebrechorkos et al., Nature, 2025 & ``area-weighted grid cells were not used when calculating study-wide averages, an error that affected numerical values in the main text'' \\
Islam et al., PLOS ONE, 2019 & ``an error in the p-value calculation caused by a coding error'' \\
Zhang et al., PLOS ONE, 2025 & a formula ``contains an incorrect minus sign, which should be a plus sign'' \\
Okray et al., Nature, 2023 & the authors found ``errors in the pipeline originally used for data analysis'' \\
Hewitson et al., PLOS ONE, 2021 & the correlation analysis ``included typically developing (TD) participants in error'' \\
Rahman et al., PLOS ONE, 2021 & ``the PSS mean reported in [1] is consistent with incorrect coding of these items'' Also states that data, images or records could not be produced. \\
Huang et al., Lancet, 2021 & an expression of concern was published ``because of the substantive data correction request from the authors following comments from a reader'' \\
Reinhart and Rogoff, American Economic Review, 2010 & ``coding errors, selective exclusion of available data, and unconventional weighting of summary statistics lead to serious errors'' \\
\midrule
\multicolumn{2}{@{}l}{\textbf{Data fabricated or falsified}} \\
Potti et al., Nature Medicine, 2006 & ``Respondent altered data sets to improve the accuracy of predictors for response to treatments''; ``the results incorporate several simple errors that may be putting patients at risk''; ``a corruption of several validation data sets precludes conclusions regarding these signatures'' \\
Potts Kant et al., Pediatric Research, 2013 & falsified data by ``changing numerical values in the datasets generated by the apparatus'' \\
Savine and Braver, Cognitive, Affective and Behavioral Neuroscience, 2012 & ``Respondent exaggerated (1) the effect of incentive context on response times and error rates in Table 1'' \\
Skau et al., Cell, 2016 & falsification by ``selectively reporting by inappropriate inclusion/omission or alteration of data points in ten (10) figures'' \\
Hauser et al., Cognition, 2002 & ``Respondent published fabricated data in Figure 2 of the paper'' \\
Linn et al., Cellular and Molecular Biology, 2003 & ``he falsified the values in Tables 2-5 by altering the values that he had obtained from another author's manuscript'' \\
Thiruchelvam et al. (publication to be identified from the finding) & misconduct by ``falsifying and fabricating cell count data that she claimed to have obtained through stereological methods'' \\
Taylor and Gercel-Taylor, Gynecologic Oncology, 2008 & two investigations ``found evidence of falsification, fabrication, and plagiarism of data'' \\
\midrule
\multicolumn{2}{@{}l}{\textbf{Image manipulation or duplication}} \\
Jiang et al., Nature, 2002 & ``The original images for Figures 6a, 6e and 6m could not be retrieved by the authors'' \\
Lesné et al., Nature, 2006 & the figures ``show signs of excessive manipulation, including splicing, duplication and the use of an eraser tool. The data cannot be verified from the records.'' \\
Zou et al., Nature, 2001 & falsified figures by ``manipulating the images to alter the number and location of positively stained cells'' \\
Liu et al., PLOS ONE, 2014 & ``the data underlying the panels of concern are no longer available'' \\
\midrule
\multicolumn{2}{@{}l}{\textbf{Study conducted differently from its description}} \\
Wakefield et al., Lancet, 1998 & the paper's claims ``that children were `consecutively referred' and that investigations were `approved' by the local ethics committee have been proven to be false'' \\
Estruch et al., New England Journal of Medicine, 2013 & ``Because of irregularities in the randomization procedures, we wish to retract the following article'' \\
Gautret et al., International Journal of Antimicrobial Agents, 2020 & ``The journal has not been able to establish whether all patients could have entered into the study in time'' Also states that data, images or records could not be produced. \\
LaCour and Green, Science, 2014 & ``Survey incentives were misrepresented.''; ``LaCour has not produced the original survey data from which someone else could independently confirm the validity of the reported findings'' \\
\midrule
\multicolumn{2}{@{}l}{\textbf{None of the above}} \\
Fukuhara et al., Science, 2005 & ``we have found that not all preparations of visfatin bind to and activate the insulin receptor'' \\
\end{longtable}
}


\paragraph{Limitations.} The 34 cases are illustrations, at least one per failure mode.
Table~\ref{tab:fc-dataset} reports how often each reason for retraction occurs among 43,364
retractions in the Retraction Watch dataset. The reason tags reflect the dataset curators' reading of each notice, and
27.6\% of the retractions analyzed are tagged as having a notice with limited or no information.
In the full dataset, 1,725 of the 2,464 retractions tagged for unavailable data or images also
have an image-related tag, so most concern missing original images. The analysis shows which
failure modes the records of Reproducible Science would expose; it does not show that any of these
cases would have been corrected sooner.
